\chapter{Literature Review}
\label{ch:literature}

The integration of artificial intelligence into document processing 
and diagram recognition has seen rapid development in recent years, 
particularly in enhancing automated understanding of visual 
information. Researchers and industry practitioners have explored 
various approaches to digitizing diagrams, recognizing shapes, 
extracting text, and understanding structural relationships. This 
chapter critically reviews recent literature related to computer 
vision for diagram recognition, optical character recognition for 
handwritten text, AI-powered code and flowchart generation, 
interactive diagramming tools, and optical graph recognition 
techniques.

The technologies employed in diagram digitization include 
convolutional neural networks (CNNs) for shape detection, optical 
character recognition (OCR) engines for text extraction, large 
language models (LLMs) for understanding diagram context, and 
interactive web frameworks for diagram editing. However, the 
literature consistently identifies several unresolved limitations, 
including accuracy challenges with hand-drawn diagrams, difficulty 
in detecting connection paths, poor handling of overlapping elements, 
and the absence of integrated solutions that combine automated 
recognition with interactive editing capabilities. These gaps 
directly motivate the design of the Automated Diagram Digitization 
system.

%---------------------------------------------------------------
\section{Computer Vision for Diagram Recognition}
\label{sec:cv-recognition}
%---------------------------------------------------------------

Vishnu Vardhan Reddy and Rajeswari present a 
unified framework that merges optical character recognition with 
generative intelligence for automated document processing. The 
framework incorporates preprocessing techniques such as noise 
reduction and layout analysis to enhance recognition accuracy, and 
employs advanced language processing to produce semantically rich, 
structured outputs from raw visual content. The authors report 
measurable improvements in extraction accuracy for structured 
documents when preprocessing pipelines are applied prior to 
recognition.~\cite{ocr-review}

However, a critical limitation of this framework is its primary 
focus on textual documents rather than diagrammatic structures. 
The system does not address the recognition of geometric shapes, 
directional arrows, or node-to-node connections, all of which are 
essential requirements for flowchart digitization. Furthermore, 
the approach does not provide an interactive editing environment 
following recognition, leaving the user with static structured 
output rather than a manipulable diagram.

This work is directly relevant to the Automated Diagram Digitization 
project as it validates the use of preprocessing pipelines and 
generative AI for structured information extraction. The preprocessing 
strategies described -- particularly noise reduction and layout 
normalization -- inform the image resizing and thumbnail generation 
steps applied before submitting images to the Gemini AI 
API.~\cite{diagram-recognition}

%---------------------------------------------------------------
\section{Handwriting Recognition in Document Processing}
\label{sec:handwriting}
%---------------------------------------------------------------

Research in deep learning-based handwriting recognition, as surveyed 
by Szeliski, demonstrates that CNNs combined with 
sequence modeling architectures such as recurrent neural networks 
(RNNs) and connectionist temporal classification (CTC) can 
effectively transcribe handwritten text with high accuracy on 
standard benchmarks. Attention-based transformer architectures have 
further improved performance by capturing long-range character 
dependencies that RNN-based models struggle to model.~\cite{szeliski}

Despite this progress, handwriting recognition systems exhibit 
significant limitations when applied to diagram labels. They struggle 
with varied handwriting styles and cursive text, are sensitive to 
image quality and lighting inconsistencies, and require large, 
diverse training datasets to generalize well. Critically, these 
systems treat text recognition as a linear sequence problem and 
cannot preserve the spatial formatting -- such as multi-line labels 
within shapes -- that is essential for accurate flowchart 
reconstruction. The challenge of preserving line breaks within 
diagram node labels is a specific gap that existing handwriting 
recognition pipelines do not address.

This limitation is directly relevant to the present project. The 
Automated Diagram Digitization system addresses it through custom 
prompt engineering that instructs the Gemini AI model to use 
\texttt{\textbackslash n} characters explicitly for line breaks 
within extracted node labels, preserving the original text 
formatting from the source diagram.~\cite{opencv}

%---------------------------------------------------------------
\section{AI-Powered Code and Flowchart Generation}
\label{sec:ai-generation}
%---------------------------------------------------------------

The application of large language models (LLMs) to automated 
flowchart generation from natural language descriptions has been 
explored through models such as GPT and Google's 
Gemini. These systems accept natural language 
algorithm descriptions as input, decompose them into logical steps, 
map each step to an appropriate flowchart symbol, and generate 
structured outputs -- typically in JSON or XML -- that represent 
the diagram's nodes and edges. Such tools aim to bridge the gap 
between conceptual algorithm design and its visual representation, 
making diagram creation accessible to users with limited diagramming 
experience.~\cite{gemini}

However, the literature identifies several persistent limitations in 
LLM-based generation. Generated outputs exhibit 
inconsistency across repeated invocations for the same input, 
particularly in node naming conventions and connection routing. 
Models struggle with complex nested logic structures such as 
deeply nested conditionals and loop-back edges. Additionally, 
post-generation validation is typically required, as models 
occasionally produce structurally invalid JSON or omit required 
fields such as edge identifiers. These limitations necessitate 
robust error handling on the application layer.~\cite{summarization}

This work is central to the Automated Diagram Digitization project. 
The text-to-flowchart feature directly implements this LLM-based 
generation approach using the Gemini API, and the limitations 
identified here directly informed the design of the backend's 
regex-based JSON extraction fallback and the edge ID auto-generation 
mechanism described in Chapter~4.

%---------------------------------------------------------------
\section{Interactive Diagramming Tools and Libraries}
\label{sec:diagramming-tools}
%---------------------------------------------------------------

Web-based diagramming libraries such as React Flow, GoJS, and 
mxGraph have been studied for their capabilities in supporting 
interactive, editable diagram 
interfaces. These libraries provide 
foundational building blocks including drag-and-drop node placement, 
edge connection handling, zoom and pan navigation, and export 
functionality. React Flow, in particular, offers a declarative 
React-based API for rendering custom node types and edges, making 
it well-suited for domain-specific diagramming applications.~\cite{reactflow, react}

The literature notes that while these libraries accelerate 
development through extensive APIs and community support, they 
exhibit performance degradation when rendering very large diagrams 
containing hundreds of nodes simultaneously. Additionally, 
implementing custom behaviors such as shape-specific resizing 
constraints, inline edge label editing, and snapshot-based undo/redo 
requires significant application-level development beyond what the 
libraries provide natively. Export format support also varies, with 
most libraries requiring additional third-party integration for PNG 
and SVG generation.

These findings directly shaped the architectural decisions of the 
Automated Diagram Digitization editor. React Flow was selected as 
the rendering engine due to its declarative component model and 
custom node support. The limitations identified -- particularly 
the absence of native undo/redo and export -- were addressed 
through the snapshot-based history system and the integration of 
the \texttt{html-to-image} library described in Chapter~5.

%---------------------------------------------------------------
\section{Optical Graph Recognition Techniques}
\label{sec:ogr}
%---------------------------------------------------------------

Kim and Lee present a dedicated study 
on optical graph recognition (OGR) -- the automated conversion of 
diagram images into machine-readable graph representations. Their 
system employs Hough transforms for line and arrow detection, 
contour analysis for shape boundary identification, and 
graph-theoretic methods for inferring node connectivity. The 
authors report successful recognition of printed flowchart diagrams 
with regular shape boundaries under controlled imaging conditions.~\cite{flowchart-analysis}

However, the approach demonstrates significant limitations when 
applied to hand-drawn or photographed diagrams. The reliance on 
traditional computer vision heuristics makes the system sensitive 
to variations in line thickness, shape irregularity, and arrow 
style. The study reports reduced accuracy for overlapping elements 
and diagrams with non-standard arrow styles. Critically, the OGR 
system produces only a static machine-readable output and provides 
no interactive editing environment, requiring users to transfer the 
recognized structure into a separate diagramming tool for further 
modification.

This work is directly relevant to the Automated Diagram Digitization 
project as it establishes the core technical challenge of automated 
shape and connection recognition. The present system addresses the 
limitations of traditional OGR by replacing heuristic computer 
vision methods with the multimodal understanding capabilities of 
the Gemini AI model, which can interpret shape semantics 
contextually rather than through rigid geometric rules. This shift 
from rule-based to AI-based recognition is the primary technical 
contribution of the proposed system.

%---------------------------------------------------------------
% Transition Paragraph (C2)
%---------------------------------------------------------------

The works reviewed above collectively highlight the significant 
progress made in diagram recognition, handwriting analysis, 
AI-assisted flowchart generation, and interactive diagramming 
tools. However, a common limitation across all reviewed approaches 
is the absence of an integrated solution that combines automated 
recognition with a feature-rich interactive editor. Systems that 
perform recognition do not provide editing, and systems that 
provide editing require entirely manual input. The proposed 
Automated Diagram Digitization system directly addresses this gap 
by combining Gemini AI-powered recognition with a React Flow-based 
interactive editor within a single, unified web platform.

%---------------------------------------------------------------
\section{Conclusion}
\label{sec:lit-conclusion}
%---------------------------------------------------------------

The reviewed literature establishes the technical foundations 
relevant to the Automated Diagram Digitization project while also 
revealing clear gaps that the proposed system is designed to 
address. The proposed system overcomes these limitations by 
providing the following capabilities:

\begin{enumerate}

    \item A vision-based recognition system that automatically 
    identifies flowchart shapes, text, and connections from uploaded 
    images using Google's Gemini AI, replacing brittle heuristic 
    computer vision methods with contextual multimodal understanding.

    \item Support for multiple input methods including image upload, 
    natural language descriptions for text-to-flowchart generation, 
    and JSON import for existing flow data -- a combination not 
    offered by any single system reviewed in this chapter.

    \item An interactive editor with comprehensive features including 
    drag-and-drop from shape palette, undo/redo with snapshot 
    history, auto-layout for clean alignment, and edge labeling via 
    double-click -- addressing the post-recognition editing gap 
    identified in OGR and OCR literature.

    \item Text customization options including bold, italic, and font 
    size adjustments, with line break preservation using 
    \texttt{\textbackslash n} characters -- directly addressing the 
    formatting loss problem identified in handwriting recognition 
    research.

    \item Multiple export formats including PNG for sharing, SVG for 
    vector graphics, and JSON for data portability -- extending 
    beyond the limited export options noted in diagramming library 
    literature.

    \item Theme switching (dark/light mode) with persistent user 
    preference and persistent storage using PostgreSQL with 
    UUID-based identification.

\end{enumerate}

Unlike the reviewed systems, which address recognition or editing 
in isolation, the Automated Diagram Digitization platform integrates 
both capabilities within a single, cohesive workflow. This 
integration represents the primary contribution of the project to 
the field of intelligent document processing. \textit{Refer Table~\ref{tab:gap_analysis}: Research Gaps Addressed by the Proposed System.}
\vspace{1em}
% Gap Analysis Table (D2)
\begin{table}[h]
\centering
\renewcommand{\arraystretch}{1.4}
\begin{tabular}{|p{8cm}|p{8cm}|}
\hline
\textbf{Gap Identified in Literature} & 
\textbf{How the Proposed System Addresses It} \\
\hline
No integrated solution combining automated recognition with 
interactive editing & 
Combines Gemini AI recognition with a React Flow interactive 
editor in a single platform \\
\hline
Poor preservation of text formatting (line breaks) in diagram 
labels & 
Uses Gemini's multimodal prompt engineering with explicit 
\texttt{\textbackslash n} line break instructions \\
\hline
Limited or single input methods in existing tools & 
Provides three input modes: image upload, text-to-flowchart 
generation, and JSON import \\
\hline
Restricted export options in automated recognition systems & 
Supports PNG, SVG, and JSON export formats for diverse 
use cases \\
\hline
Heuristic OGR methods fail on hand-drawn or irregular diagrams & 
Replaces rule-based vision with contextual multimodal AI 
understanding via Gemini API \\
\hline
\end{tabular}
\caption{Research Gaps Addressed by the Proposed System}
\label{tab:gap_analysis}
\end{table}